% !TEX root = ../report.tex
\chapter{Аналитическая часть}
\section{Постановка задачи}
Дан массив целых чисел $a$ длины $N$ и искомое значение $x$. Нужно вернуть индекс элемента, равного $x$, или $-1$, если такого элемента нет. Индексы начинаются с нуля. При повторяющихся значениях допустим индекс любого равного элемента. Для пустого массива результат равен $-1$.

Линейный алгоритм получает исходный массив. Бинарные алгоритмы получают его копию, отсортированную по неубыванию. Линейный поиск возвращает индекс в исходном массиве, а бинарные --- в отсортированной копии.

\section{Линейный поиск}
Линейный поиск последовательно проверяет элементы с индексами от нуля до $N-1$ и завершает работу при первом равенстве. Если проверены все элементы и равенство не обнаружено, возвращается $-1$.

Линейный поиск не требует сортировки массива~\cite{cormen}.

\section{Классический нерекурсивный бинарный поиск}
Область поиска задаётся полуинтервалом $[L,R)$. На каждом шаге вычисляется индекс среднего элемента. Если этот элемент меньше $x$, левая граница переносится за него. Если он больше $x$, правая граница становится равной его индексу. При равенстве возвращается индекс среднего элемента.

Благодаря упорядоченности массива исключённая часть не содержит искомого значения. После каждого неудачного шага остаётся не более половины элементов предыдущего участка с округлением вниз.

\section{Бинарный поиск без ветвления}
Второй нерекурсивный алгоритм взят из статьи М. Скарупке~\cite{skarupke}. Поиск выполняется с шагом \texttt{step}, который последовательно уменьшается вдвое. Для выбора начального шага используются две функции: \texttt{bit\_floor(N)} возвращает наибольшую степень двойки, не превышающую $N$, а \texttt{bit\_ceil(N)} --- наименьшую степень двойки, не меньшую $N$. Например, для $N=10$ это 8 и 16 соответственно.

Сначала \texttt{begin} равен нулю, а \texttt{step} равен \texttt{bit\_floor(N)}. Если размер массива не является степенью двойки, алгоритм сравнивает элемент с индексом \texttt{step} с искомым значением. Если этот элемент меньше искомого, поиск переносится в конец массива. Для оставшейся части заново выбирается шаг с помощью \texttt{bit\_ceil}. Если за проверенным элементом больше ничего нет, возвращается $-1$.

В основном цикле шаг каждый раз делится на два. Проверяется элемент с индексом \texttt{begin + step}. Если он меньше искомого значения, \texttt{begin} увеличивается на \texttt{step}; иначе \texttt{begin} остаётся прежним. Цикл заканчивается, когда шаг становится равным нулю.

После цикла выполняется ещё одно сравнение элемента с индексом \texttt{begin} с искомым значением. Оно определяет позицию первого элемента, который не меньше искомого. В статье алгоритм возвращает эту позицию, даже если самого значения в массиве нет. В данной программе дополнительно проверяется равенство: если значение найдено, возвращается его индекс, иначе --- $-1$.

Название «без ветвления» относится к машинному коду основного цикла. В исходном алгоритме есть условие \texttt{if}, которое GCC может преобразовать в инструкцию условного перемещения. Поэтому отсутствие условного перехода зависит от компилятора и параметров сборки; оно не означает отсутствие \texttt{if} в исходном коде.

\section{Рекурсивный бинарный поиск}
На каждом вызове проверяется размер текущего массива. При нулевом размере возвращается $-1$. Иначе сначала проверяется равенство среднего элемента искомому значению. Если значения различаются, выбирается левая часть массива, когда искомое значение меньше среднего элемента, и правая --- когда больше.

При поиске в правой части функция возвращает индекс относительно этой части. Чтобы получить индекс в исходном массиве, к нему прибавляется число элементов, оставшихся слева. При поиске в левой части индекс не меняется. Каждый следующий вызов получает меньше элементов, поэтому поиск заканчивается. В каждом вызове выбирается только одна часть массива, и рекурсивные вызовы образуют цепочку.

\section*{Вывод}
Сформулирована задача поиска и рассмотрены четыре алгоритма её решения. Линейный поиск применим к неупорядоченным данным, а бинарные алгоритмы требуют предварительного упорядочивания. Два нерекурсивных бинарных алгоритма различаются порядком сравнений и возможностью раннего завершения.
